% !TEX program = xelatex
\documentclass[12pt,a4paper]{report}

\usepackage{graphicx}
\usepackage{geometry}
\usepackage{float}
\usepackage{xcolor}
\usepackage{titlesec}
\usepackage{xepersian}

\geometry{top=3cm, bottom=3cm, left=2.5cm, right=3cm}

% تنظیم فاصله بعد از عنوان‌ها
\titlespacing*{\section}{0pt}{2ex plus 1ex minus .2ex}{1.5ex plus .2ex}
\titlespacing*{\subsection}{0pt}{1.8ex plus 1ex minus .2ex}{1ex plus .2ex}

% فونت فارسی ۱۴ و انگلیسی ۱۲
\settextfont[
Path = C:/USERS/SAZGAR/APPDATA/LOCAL/MICROSOFT/WINDOWS/FONTS/,
BoldFont = IRNAZANIN.TTF,
Scale = 1.166
]{IRNAZANIN.TTF}
\setlatintextfont[Scale=1]{Times New Roman}

\linespread{1.5}

\begin{document}
	\sloppy
	\chapter{رویکرد \lr{Zero-shot}}
	
	در گام نخست این پژوهش، عملکرد ذاتی \lr{(Baseline Performance)} چهارده مدل زبانی بزرگ در پاسخ به \lr{120} سوال از آزمون وکالت سال ۱۴۰۳ را مورد ارزیابی قرار داد. این آزمایش در حالت \lr{Zero-Shot}، یعنی بدون ارائه هرگونه مثال یا راهنمایی \lr{(In-context Learning)}، انجام شد تا دانش درونی و توانایی استدلال خام مدل‌ها بدون دسترسی به منابع خارجی سنجیده شود[1].
	
	شکل \ref{fig:zero-shot-acc} مقایسه دقت مدل‌ها را با تفکیک دو گروه متن‌باز \lr{(Open Source)} و تجاری \lr{(Closed Source)} نمایش می‌دهد. نتایج نشان‌دهنده یک عملکرد کلی متوسط است که در آن هیچ مدلی نتوانست به دقت چشمگیری دست یابد. مدل متن‌باز \lr{Qwen3-235} با کسب دقت \lr{\%52.5} به طور غیرمنتظره‌ای در جایگاه نخست قرار گرفت و فرضیه برتری مطلق مدل‌های تجاری را به چالش کشید. پس از آن، مدل‌های \lr{DeepSeek-R1-T2-Chimera} (با دقت \lr{\%49.2}) و \lr{Gemini 2.5 Flash} (با دقت \lr{\%48.3}) به ترتیب در رتبه‌های دوم و سوم قرار گرفتند. در مقابل، مدل‌های خانواده \lr{GPT-OSS} با دقت تنها \lr{\%32.5} در انتهای جدول جای گرفتند که بیانگر کمبود داده‌های حقوقی فارسی در فرآیند پیش‌آموزش آن‌هاست.
	
	\begin{figure}[H]
		\centering
		\includegraphics[width=0.8\textwidth]{F:/Thesis/project/1-BaselineTest/Models Responses/zero-shot/++models analysis++/zero-shot-acc.png}
		\caption{مقایسه دقت مدل‌ها در حالت \lr{Zero-shot} با تفکیک مدل‌های متن‌باز و تجاری.}
		\label{fig:zero-shot-acc}
	\end{figure}
	
	برای کاربردهای عملی، سرعت پاسخ‌دهی یک عامل تعیین‌کننده است. شکل \ref{fig:latency} زمان پاسخ‌دهی \lr{(Latency)} مدل‌ها را به میلی‌ثانیه نشان می‌دهد. یک شکاف عملکردی عظیم در این بخش مشاهده می‌شود؛ مدل \lr{DeepSeek-R1-T2-Chimera} با وجود دقت بالا، تأخیری معادل \lr{18.5} ثانیه داشت که آن را برای کاربردهای بلادرنگ نامناسب می‌سازد. در مقابل، خانواده مدل‌هایی مانند \lr{Gemini 2.5} و \lr{Llama 4} با سرعتی نزدیک به \lr{600} تا \lr{700} میلی‌ثانیه، کارایی زمانی بسیار بهتری را به نمایش گذاشتند. این اختلاف چشمگیر در سرعت، احتمالاً به دلیل تفاوت در پیچیدگی معماری و بهینه‌سازی‌های استنتاجی \lr{(Inference Optimization)} است[2].
	
	\begin{figure}[H]
		\centering
		\includegraphics[width=0.8\textwidth]{F:/Thesis/project/1-BaselineTest/Models Responses/zero-shot/++models analysis++/zero shot latency.png}
		\caption{میانگین زمان پاسخ‌دهی \lr{(Latency)} مدل‌ها به میلی‌ثانیه.}
		\label{fig:latency}
	\end{figure}
	
	شکل \ref{fig:speed-acc} رابطه بین دو معیار کلیدی سرعت و دقت را در یک نمودار پراکندگی به تصویر می‌کشد. این نمودار به وضوح نشان می‌دهد که مدل \lr{Qwen3-235} (نقطه بالا و سمت چپ) بهترین تعادل \lr{(Trade-off)} را میان این دو معیار برقرار کرده است. خوشه‌ای از مدل‌ها مانند \lr{Gemini 2.5 Flash} نیز در ناحیه «دقت بالا و سرعت بالا» قرار دارند که آن‌ها را به گزینه‌هایی کارآمد برای پیاده‌سازی تبدیل می‌کند.
	
	\begin{figure}[H]
		\centering
		\includegraphics[width=0.8\textwidth]{F:/Thesis/project/1-BaselineTest/Models Responses/zero-shot/++models analysis++/speed vs acc.png}
		\caption{نمودار پراکندگی دقت (محور عمودی) در برابر تأخیر (محور افقی)؛ رابطه بین تأخیر پاسخ‌دهی و دقت مدل‌ها.}
		\label{fig:speed-acc}
	\end{figure}
	
	در نهایت، شکل \ref{fig:confidence} به بررسی کالیبراسیون اطمینان مدل‌ها می‌پردازد. یکی از چالش‌های جدی در استفاده از مدل‌های زبانی، پدیده «توهم با اطمینان بالا» \lr{(High-confidence Hallucination)} است[3]. مقایسه این نمودار با نمودار دقت (شکل \ref{fig:zero-shot-acc}) یک عدم تطابق نگران‌کننده را آشکار می‌سازد. برای مثال، مدل \lr{Llama 4 Maverick} با وجود دقت متوسط (برابر با \lr{\%42.5})، بالاترین امتیاز اطمینان (\lr{4.21}) را به خود اختصاص داده است. این رفتار که به عنوان «اعتماد به نفس کاذب» شناخته می‌شود، در حوزه‌های حساس مانند حقوق بسیار خطرناک است. در مقابل، مدل \lr{DeepSeek-R1-T2} با ثبت کمترین امتیاز اطمینان (\lr{3.56})، دقت بالایی را کسب کرد که نشان‌دهنده کالیبراسیون بهتر و توانایی مدل در تشخیص عدم قطعیت‌های خود است[4].
	
	\begin{figure}[H]
		\centering
		\includegraphics[width=0.8\textwidth]{F:/Thesis/project/1-BaselineTest/Models Responses/zero-shot/++models analysis++/zero shot confidene.png}
		\caption{میانگین امتیاز اطمینان \lr{(Confidence Score)} مدل‌ها.}
		\label{fig:confidence}
	\end{figure}
	
	در مجموع، نتایج فاز \lr{Zero-shot} نشان داد که دانش پارامتریک مدل‌ها به تنهایی برای پاسخگویی قابل اتکا به سوالات پیچیده حقوقی کافی نیست و این امر ضرورت بررسی روش‌های پیشرفته‌تری مانند \lr{Chain-of-Thought} و \lr{Retrieval-Augmented Generation (RAG)} را که در فازهای بعدی این پژوهش به آن‌ها پرداخته می‌شود، بیش از پیش آشکار می‌سازد.
	
	
\section{رویکرد \lr{Few-shot}}

در گام دوم ارزیابی، تأثیر یادگیری در لحظه \lr{(In-context Learning)} بر توانایی استدلال حقوقی مدل‌ها مورد بررسی قرار گرفت. در رویکرد \lr{Few-shot}، به هر مدل \lr{6} نمونه سوال حل‌شده (شامل سوال، گزینه صحیح و استدلال مختصر) به عنوان ورودی \lr{(Prompt)} داده شد تا الگوی پاسخگویی را فراگیرد[1]. هدف از این آزمایش، سنجش این فرضیه بود که آیا ارائه نمونه‌های مشابه می‌تواند شکاف دانش مشاهده شده در حالت \lr{Zero-shot} را کاهش دهد یا خیر.

شکل \ref{fig:few-shot-acc} مقایسه دقت مدل‌ها را در حالت \lr{Few-shot} نشان می‌دهد. برخلاف حالت \lr{Zero-shot} که مدل‌های متن‌باز پیشتاز بودند، در اینجا مدل تجاری \lr{Gemini 2.5 Flash} با جهشی قابل توجه به دقت \lr{\%55.0} دست یافت و رتبه نخست را کسب کرد. مدل \lr{Grok 4 Fast} نیز با دقت \lr{\%54.2} در جایگاه دوم قرار گرفت. این تغییر در رتبه‌بندی نشان می‌دهد که مدل‌های تجاری بزرگ‌تر \lr{(Closed Source)} توانایی بالاتری در بهره‌گیری از نمونه‌های آموزشی \lr{(Few-shot prompt)} دارند[5]. در میان مدل‌های متن‌باز، \lr{Llama 4 Maverick} و \lr{Qwen3-235} با دقت برابر \lr{\%45.8} بهترین عملکرد را داشتند، اما همچنان نسبت به رقبای تجاری خود عقب ماندند. مدل‌های خانواده \lr{GPT-OSS} همچنان با دقت‌های زیر \lr{\%40}، عملکرد ضعیف خود را تکرار کردند.

\begin{figure}[H]
	\centering
	\includegraphics[width=0.8\textwidth]{F:/Thesis/project/1-BaselineTest/Models Responses/few-shot/++models analysis++/few_shot_accuracy.png}
	\caption{مقایسه دقت مدل‌ها در حالت \lr{Few-shot} با تفکیک مدل‌های متن‌باز و تجاری.}
	\label{fig:few-shot-acc}
\end{figure}

از نظر کارایی زمانی، نتایج نمایش داده شده در شکل \ref{fig:few-latency} همچنان برتری مطلق مدل‌های سبک‌تر را تایید می‌کند. مدل‌های خانواده \lr{Gemini 2.5} با میانگین زمان پاسخگویی حدود \lr{620} میلی‌ثانیه، سریع‌ترین عملکرد را ثبت کردند. در نقطه مقابل، مدل \lr{DeepSeek-R1-T2-Chimera} با تأخیری بیش از \lr{14} ثانیه، کندترین مدل ارزیابی شده بود. نکته قابل تامل این است که افزودن نمونه‌های \lr{Few-shot} به ورودی مدل (که طول توکن‌های ورودی را افزایش می‌دهد)، تأثیر منفی معناداری بر سرعت مدل‌های بهینه شده (مانند \lr{Gemini} و \lr{Llama}) نداشته است، اما مدل‌های سنگین‌تر را کندتر کرده است[6].

\begin{figure}[H]
	\centering
	\includegraphics[width=0.8\textwidth]{F:/Thesis/project/1-BaselineTest/Models Responses/few-shot/++models analysis++/few shot latency.png}
	\caption{میانگین زمان پاسخ‌دهی \lr{(Latency)} مدل‌ها در حالت \lr{Few-shot}.}
	\label{fig:few-latency}
\end{figure}

شکل \ref{fig:few-speed-acc} تحلیل جامع‌تری از رابطه میان سرعت و دقت ارائه می‌دهد. در این نمودار پراکندگی، مدل \lr{Gemini 2.5 Flash} (نقطه زرد رنگ در بالا سمت چپ) به عنوان «نقطه طلایی» \lr{(Sweet Spot)} شناخته می‌شود؛ زیرا بالاترین دقت را با یکی از کمترین تأخیرها ترکیب کرده است. این در حالی است که مدل \lr{DeepSeek-R1-T2} (نقطه قرمز سمت راست) علیرغم دقت قابل قبول، به دلیل تأخیر بسیار بالا، از ناحیه کارآمدی خارج شده است. این نمودار به وضوح نشان می‌دهد که برای کاربردهای عملیاتی حقوقی، مدل‌هایی که در نیمه چپ و بالای نمودار قرار دارند (مانند \lr{Gemini} و \lr{Grok})، گزینه‌های برتر هستند.

\begin{figure}[H]
	\centering
	\includegraphics[width=0.8\textwidth]{F:/Thesis/project/1-BaselineTest/Models Responses/few-shot/++models analysis++/speed vs acc.png}
	\caption{نمودار تعامل سرعت و دقت \lr{(Speed vs Accuracy Trade-off)} در حالت \lr{Few-shot}.}
	\label{fig:few-speed-acc}
\end{figure}

نهایتاً، شکل \ref{fig:few-confidence} وضعیت کالیبراسیون اطمینان مدل‌ها را بررسی می‌کند. یک یافته مهم در این بخش، کاهش کلی میانگین اطمینان در برخی مدل‌های دقیق‌تر است. برای مثال، مدل \lr{DeepSeek-R1-T2} میانگین اطمینان خود را به \lr{3.24} کاهش داده است که نشان‌دهنده احتیاط بیشتر مدل در مواجهه با پرامپت‌های طولانی‌تر است. در مقابل، مدل \lr{GPT-4.1 Mini} با وجود دقت پایین (\lr{\%41.7})، به طور نگران‌کننده‌ای بالاترین سطح اطمینان (\lr{4.53}) را گزارش کرده است. این «اطمینان کاذب» \lr{(Overconfidence)} در مدل‌های ضعیف‌تر، ریسک استفاده از آن‌ها را در سامانه‌های حساس حقوقی افزایش می‌دهد، زیرا توانایی تشخیص مرزهای دانش خود را ندارند[7].

\begin{figure}[H]
	\centering
	\includegraphics[width=0.8\textwidth]{F:/Thesis/project/1-BaselineTest/Models Responses/few-shot/++models analysis++/few shot confidene.png}
	\caption{میانگین امتیاز اطمینان \lr{(Confidence Score)} مدل‌ها در حالت \lr{Few-shot}.}
	\label{fig:few-confidence}
\end{figure}

در مجموع، گذار از \lr{Zero-shot} به \lr{Few-shot} منجر به بهبود عملکرد مدل‌های تجاری (به ویژه \lr{Gemini} و \lr{Grok}) شد، اما مدل‌های متن‌باز تغییر چشمگیری را تجربه نکردند. این امر حاکی از آن است که معماری مدل‌های تجاری احتمالاً برای یادگیری در متن \lr{(In-context Learning)} بهینه‌تر شده است. با این حال، حتی بهترین دقت کسب شده (\lr{\%55}) هنوز برای یک سامانه مشاور حقوقی قابل اتکا کافی نیست و نیاز به روش‌های استدلالی عمیق‌تر (مانند \lr{Chain-of-Thought}) را برجسته می‌کند[8].

\section{‌تأثیر یادگیری در متن \lr{(Zero-shot vs. Few-shot)}}

برای درک دقیق‌تر تأثیر افزودن نمونه‌های آموزشی به ورودی مدل‌ها، مقایسه‌ای مستقیم بین عملکرد \lr{Zero-shot} و \lr{Few-shot} انجام شد.

شکل \ref{fig:grouped-bar} مقایسه پهلو-به-پهلو \lr{(Side-by-Side)} دقت مدل‌ها در دو حالت را نشان می‌دهد. همان‌طور که مشاهده می‌شود، پاسخ مدل‌ها به تکنیک \lr{Few-shot} یکسان نبوده است. مدل‌هایی مانند \lr{Grok 4 Fast} و \lr{Gemini 2.5 Flash} جهش مثبتی را تجربه کرده‌اند که نشان‌دهنده توانایی بالای آن‌ها در انطباق با الگوهای جدید است. در مقابل، برخی مدل‌های قدرتمند در حالت \lr{Zero-shot} مانند \lr{Qwen3-235}، در مواجهه با نمونه‌های \lr{Few-shot} دچار افت عملکرد شده‌اند.

\begin{figure}[H]
	\centering
	\includegraphics[width=0.85\textwidth]{F:/Thesis/project/1-BaselineTest/Models Responses/few-shot/++models analysis++/Grouped Bar.png}
	\caption{مقایسه دقت مدل‌ها در دو حالت \lr{Zero-shot} (آبی روشن) و \lr{Few-shot} (سرمه‌ای).}
	\label{fig:grouped-bar}
\end{figure}

برای کمی‌سازی این تغییرات، شکل \ref{fig:delta} میزان تغییر دقت \lr{(Delta Accuracy)} را به درصد نشان می‌دهد. نتایج شگفت‌انگیز است:

\textbf{بزرگترین برندگان:} مدل \lr{Grok 4 Fast} با \lr{+\%10} افزایش دقت، بیشترین بهره را از نمونه‌های \lr{Few-shot} برده است. پس از آن، مدل‌های خانواده \lr{GPT-OSS} و \lr{Gemini 2.5 Flash} با رشدهای \lr{\%5} تا \lr{\%6.7} قرار دارند. این رشد نشان می‌دهد که این مدل‌ها پتانسیل نهفته‌ای داشتند که با هدایت صحیح \lr{(Prompting)} فعال شده است.

\textbf{بزرگترین بازندگان:} در کمال تعجب، مدل \lr{Qwen3-235} که قهرمان حالت \lr{Zero-shot} بود، با افت \lr{-\%6.7} مواجه شد. همچنین \lr{GPT-4o-mini} (\lr{-\%5}) و \lr{DeepSeek-R1-T2} (\lr{-\%4.2}) نیز عملکرد ضعیف‌تری نشان دادند.

این پدیده که به عنوان «تداخل منفی» \lr{(Negative Transfer)} یا «حواس‌پرتی با نمونه‌ها» شناخته می‌شود، زمانی رخ می‌دهد که نمونه‌های ارائه شده \lr{(Few-shot examples)} با دانش درونی مدل تعارض پیدا می‌کنند یا باعث سردرگمی مدل در تشخیص الگوی اصلی می‌شوند[9]. این یافته مهم نشان می‌دهد که استراتژی \lr{Few-shot} همیشه راهگشا نیست و برای برخی مدل‌ها (به‌ویژه مدل‌های متن‌باز خاص)، تکیه بر دانش درونی \lr{(Zero-shot)} موثرتر است.

\begin{figure}[H]
	\centering
	\includegraphics[width=0.85\textwidth]{F:/Thesis/project/1-BaselineTest/Models Responses/few-shot/++models analysis++/delta.png}
	\caption{نمودار تغییرات دقت \lr{(Delta)} مدل‌ها پس از اعمال روش \lr{Few-shot}. مقادیر سبز نشان‌دهنده بهبود و مقادیر قرمز نشان‌دهنده افت عملکرد هستند.}
	\label{fig:delta}
\end{figure}


\section{رویکرد \lr{Zero-shot Chain of Thought (Explain-to-Predict)}}

در گام سوم ارزیابی، از روش \lr{Explain-to-Predict (E2P)} در قالب \lr{Zero-shot Chain of Thought (CoT)} استفاده شد. در این روش، از مدل خواسته می‌شود پیش از اعلام گزینه نهایی، ابتدا استدلال گام‌به‌گام خود را ارائه دهد. هدف از این آزمایش، بررسی این فرضیه است که آیا وادار کردن مدل به «تفکر صریح» (\lr{Explicit Reasoning}) می‌تواند کیفیت پاسخ‌های حقوقی را بهبود بخشد و شکاف عملکردی مشاهده شده در مراحل قبل را جبران کند یا خیر[8].

شکل \ref{fig:cot-acc} عملکرد مدل‌ها را در این سناریو نمایش می‌دهد. نتایج حاکی از یک جهش عملکردی چشمگیر در برخی مدل‌ها است. مدل \lr{Gemini 2.5 Flash} با ثبت رکورد دقت \lr{\%57.5}، بالاترین عملکرد را در میان تمامی روش‌های تست شده تا کنون (\lr{Zero-shot} و \lr{Few-shot}) به نام خود ثبت کرد. این افزایش دقت نشان می‌دهد که مکانیزم‌های استنتاجی \lr{(Reasoning Capabilities)} در مدل‌های پیشرفته تجاری، با استراتژی \lr{CoT} به خوبی فعال می‌شوند[10].

در سوی دیگر، مدل‌های \lr{Grok 4 Fast} با \lr{\%48.3} و \lr{Qwen3-235} با \lr{\%47.5} در رتبه‌های بعدی قرار گرفتند. نکته قابل توجه، بهبود نسبی عملکرد مدل‌های \lr{GPT-OSS} است که توانستند از قعر جدول فاصله بگیرند (دقت \lr{\%45} برای نسخه \lr{120})، که نشان می‌دهد ضعف دانشی آن‌ها تا حدی با استدلال منطقی قابل جبران است.

\begin{figure}[H]
	\centering
	\includegraphics[width=0.8\textwidth]{F:/Thesis/project/1-BaselineTest/Models Responses/zero-shot-cot-e2p/++models analysis++/zero_shot_cot_e2p.png}
	\caption{مقایسه دقت مدل‌ها در روش \lr{Zero-shot Chain of Thought (E2P)}.}
	\label{fig:cot-acc}
\end{figure}

فعال‌سازی زنجیره افکار هزینه سنگینی بر کارایی زمانی مدل‌ها تحمیل کرده است. شکل \ref{fig:cot-latency} نشان می‌دهد که میانگین زمان پاسخگویی در اکثر مدل‌ها به شدت افزایش یافته است. مدل \lr{DeepSeek-R1-T2-Chimera} با ثبت میانگین تأخیر خیره‌کننده \lr{20,684} میلی‌ثانیه (بیش از \lr{20} ثانیه)، عملاً برای تعاملات بلادرنگ(realtime) غیرقابل استفاده شده است[2].

با این حال، مدل \lr{Gemini 2.5 Flash Lite} شگفتی‌ساز این بخش است؛ این مدل توانسته با حفظ سرعت بسیار بالا (\lr{841} میلی‌ثانیه)، دقتی قابل قبول (\lr{\%46.7}) ارائه دهد که آن را به کارآمدترین گزینه برای کاربردهای عملی تبدیل می‌کند.

\begin{figure}[H]
	\centering
	\includegraphics[width=0.8\textwidth]{F:/Thesis/project/1-BaselineTest/Models Responses/zero-shot-cot-e2p/++models analysis++/latency.png}
	\caption{میانگین زمان پاسخ‌دهی \lr{(Latency)} مدل‌ها در فرآیند تولید استدلال و پاسخ.}
	\label{fig:cot-latency}
\end{figure}

نمودار پراکندگی در شکل \ref{fig:cot-speed-acc}، تمایز آشکاری را میان مدل‌ها ترسیم می‌کند. مدل \lr{Gemini 2.5 Flash} به تنهایی در ناحیه «دقت بالا - سرعت بالا» قرار دارد که نشان‌دهنده معماری بهینه آن برای وظایف استدلالی است. در مقابل، مدل‌های \lr{DeepSeek} و \lr{GPT-OSS-20} (نقاط سمت راست نمودار) در ناحیه «ناکارآمد» قرار گرفته‌اند، جایی که افزایش زمان پردازش منجر به بهبود متناسبی در دقت نشده است.

\begin{figure}[H]
	\centering
	\includegraphics[width=0.8\textwidth]{F:/Thesis/project/1-BaselineTest/Models Responses/zero-shot-cot-e2p/++models analysis++/speed vs acc.png}
	\caption{نمودار تعامل سرعت و دقت در حالت \lr{Zero-shot CoT}.}
	\label{fig:cot-speed-acc}
\end{figure}

شکل \ref{fig:cot-confidence} تغییرات معناداری را در الگوی اطمینان مدل‌ها نشان می‌دهد. برخلاف حالت \lr{Zero-shot} که مدل‌های ضعیف دچار اعتماد‌به‌نفس کاذب بودند، در روش \lr{CoT} شاهد یک «تعدیل اطمینان» هستیم. مدل \lr{Claude Haiku 4.5} با کاهش اطمینان خود به \lr{3.75}، رفتاری محافظه‌کارانه‌تر در پیش گرفته است. در مقابل، مدل \lr{GPT-OSS-20} همچنان با اطمینان بالا (\lr{4.35}) و دقت بسیار پایین (\lr{\%32.5})، بدترین کالیبراسیون را دارد که نشان‌دهنده ناتوانی آن در ارزیابی صحیح استدلال‌های خود است[10].

\begin{figure}[H]
	\centering
	\includegraphics[width=0.8\textwidth]{F:/Thesis/project/1-BaselineTest/Models Responses/zero-shot-cot-e2p/++models analysis++/confidence.png}
	\caption{میانگین امتیاز اطمینان مدل‌ها در روش \lr{Zero-shot CoT}.}
	\label{fig:cot-confidence}
\end{figure}

روش \lr{Zero-shot CoT} توانست سقف عملکردی مدل‌ها را به \lr{\%57.5} ارتقا دهد، اما این بهبود به قیمت افزایش چشمگیر زمان پردازش تمام شد. نتایج نشان می‌دهد که «تفکر قبل از پاسخ‌دهی» برای مدل‌های پیشرفته (مانند \lr{Gemini}) بسیار سودمند است، اما برای مدل‌های ضعیف‌تر (مانند \lr{GPT-OSS})، صرفاً منجر به تولید متن بیشتر بدون بهبود دقت می‌شود \lr{(Hallucination Amplification)}.

\section{تحلیل عمیق رفتار استدلالی و کارایی توکن}

در این بخش، به کالبدشکافی دقیق‌تر فرآیند تولید پاسخ می‌پردازیم تا درک کنیم مدل‌ها «چگونه» به پاسخ رسیده‌اند. داده‌های جدید (مانند طول استدلال، مصرف توکن و پایداری) دیدگاه‌های ارزشمندی درباره کیفیت تفکر مدل‌ها ارائه می‌دهند[12].

یکی از شاخص‌های رفتاری مهم در مدل‌های زبانی، رابطه بین «حجم متن تولیدی» و «صحت پاسخ» است. شکل \ref{fig:completion-tokens} میانگین تعداد توکن‌های تولید شده \lr{(Completion Tokens)} را برای پاسخ‌های صحیح (سبز) و نادرست (قرمز) مقایسه می‌کند. در اکثر مدل‌های مورد بررسی، به‌ویژه مدل‌های ضعیف‌تر مانند \lr{GPT-OSS-20} و \lr{Grok 4 Fast}، مشاهده می‌شود که میانگین توکن‌ها در پاسخ‌های نادرست به طور معناداری بیشتر از پاسخ‌های صحیح است. برای مثال، \lr{Grok 4 Fast} در پاسخ‌های نادرست به طور میانگین \lr{1894} توکن تولید کرده است، در حالی که برای پاسخ‌های صحیح \lr{1488} توکن مصرف کرده است. این اختلاف \lr{400} توکنی نشان می‌دهد که مدل وقتی در یافتن مسیر استدلالی صحیح ناتوان است، به «تولید بیهوده متن» (\lr{Verbose Hallucination}) روی می‌آورد و سعی می‌کند با طولانی کردن پاسخ، فقدان منطق صحیح را جبران کند[13]. در مقابل، مدل‌های کارآمدتری مانند \lr{Gemini 2.5 Flash} و \lr{Claude Haiku} رفتار متفاوتی دارند؛ تفاوت حجم توکن در پاسخ‌های درست و غلط آن‌ها بسیار ناچیز است (مثلاً در \lr{Gemini} تنها حدود \lr{2.5} توکن اختلاف وجود دارد: \lr{83.6 vs 81.2}). این «پایداری حجمی» نشان‌دهنده استراتژی استدلالی مشخص و منضبط در این مدل‌هاست.

\begin{figure}[H]
	\centering
	\includegraphics[width=0.8\textwidth]{F:/Thesis/project/1-BaselineTest/Models Responses/zero-shot-cot-e2p/++models analysis++/Completion.png}
	\caption{مقایسه میانگین توکن‌های تولید شده در پاسخ‌های صحیح (سبز) و نادرست (قرمز).}
	\label{fig:completion-tokens}
\end{figure}

شکل \ref{fig:explain-length} رابطه بین «طول توضیحات» (تعداد کلمات استدلال) و درستی پاسخ را نشان می‌دهد. یک الگوی رفتاری بسیار جالب در اینجا مشاهده می‌شود؛ استدلال‌های طولانی‌تر همیشه بهتر نیستند. مدل \lr{Claude Haiku 4.5} با تولید میانگین \lr{57} کلمه برای پاسخ‌های صحیح و \lr{56} کلمه برای غلط، پایدارترین رفتار را دارد. در مقابل، مدل‌های ضعیف‌تر مانند \lr{Grok 4 Fast} و \lr{Llama 4 Scout} در پاسخ‌های نادرست خود تمایل به تولید متن طولانی‌تری دارند (نقاط قرمز در سمت راست‌تر نمودار)، که این پدیده «پرگویی در شکست» \lr{(Verbosity in Failure)} نامیده می‌شود و نشان‌دهنده تلاش ناموفق مدل برای پوشاندن ضعف استدلالی است.

\begin{figure}[H]
	\centering
	\includegraphics[width=0.8\textwidth]{F:/Thesis/project/1-BaselineTest/Models Responses/zero-shot-cot-e2p/++models analysis++/explain len.png}
	\caption{توزیع طول توضیحات \lr{(Explain Length)} برای پاسخ‌های صحیح (سبز) و نادرست (قرمز).}
	\label{fig:explain-length}
\end{figure}

برای درک بهتر رابطه بین کارایی و دقت، شکل \ref{fig:token-efficiency} نمودار حبابی کارایی توکن را نمایش می‌دهد که در آن محور افقی «کل توکن‌های مصرفی» و محور عمودی «دقت» است. مدل‌های \lr{Gemini 2.5 Flash} و \lr{Llama 4 Scout} در گوشه بالا-چپ نمودار (خوشه کارآمد) قرار دارند؛ این یعنی آن‌ها با مصرف توکن کمتر (حدود \lr{600-700} توکن)، به دقت بالاتری دست یافته‌اند. این «ادعا در استدلال» نشانه بلوغ مدل در درک حقوقی است. در سوی دیگر، مدل \lr{Grok 4 Fast} (حباب بزرگ در سمت راست) با مصرف میانگین بیش از \lr{2500} توکن، ناکارآمدترین مدل محسوب می‌شود.

\begin{figure}[H]
	\centering
	\includegraphics[width=0.8\textwidth]{F:/Thesis/project/1-BaselineTest/Models Responses/zero-shot-cot-e2p/++models analysis++/Token Efficiency.png}
	\caption{نمودار کارایی توکن؛ مدل‌های سمت چپ و بالا کارآمدترین‌ها هستند.}
	\label{fig:token-efficiency}
\end{figure}

شکل \ref{fig:stability} شاخص «دلتا طول» (تفاوت طول استدلال در پاسخ‌های غلط نسبت به درست) را نشان می‌دهد که معیاری برای پایداری رفتاری مدل‌هاست. مدل \lr{Llama 4 Scout} وقتی اشتباه می‌کند، به طور میانگین \lr{8} کلمه بیشتر صحبت می‌کند (میله قرمز بلند)، که این یک سیگنال قوی برای تشخیص خطا است. در مقابل، مدل \lr{Claude Haiku 4.5} رفتار معکوسی دارد و در پاسخ‌های غلط کوتاه‌تر صحبت می‌کند (میله سبز منفی)، که نشان‌دهنده «تسلیم شدن سریع» در برابر ابهام است.

\begin{figure}[H]
	\centering
	\includegraphics[width=0.8\textwidth]{F:/Thesis/project/1-BaselineTest/Models Responses/zero-shot-cot-e2p/++models analysis++/Stability of Reasoning (Delta Chart).png}
	\caption{نمودار پایداری استدلال (تفاوت طول توضیحات در پاسخ‌های غلط و درست).}
	\label{fig:stability}
\end{figure}

نهایتاً، شکل \ref{fig:heatmap} نقشه حرارتی \lr{(Heatmap)} تمامی متریک‌ها را یکجا نمایش می‌دهد و به ما اجازه می‌دهد نقاط قوت و ضعف را در یک نگاه ببینیم. مدل \lr{Gemini 2.5 Flash} با رنگ‌های آبی تیره در ستون‌های \lr{Accuracy} و \lr{Efficiency}، به عنوان «قهرمان تعادل» شناخته می‌شود. در مقابل، مدل \lr{GPT-OSS-20} با رنگ‌های قرمز تیره در ستون‌های \lr{Latency} و \lr{Verbosity}، بدترین عملکرد کلی را دارد.

\begin{figure}[H]
	\centering
	\includegraphics[width=0.9\textwidth]{F:/Thesis/project/1-BaselineTest/Models Responses/zero-shot-cot-e2p/++models analysis++/Heatmap(Models × Metrics).png}
	\caption{نقشه حرارتی عملکرد مدل‌ها در تمامی متریک‌های ارزیابی \lr{CoT}.}
	\label{fig:heatmap}
\end{figure}

تحلیل‌های عمیق این بخش نشان داد که صرفاً «فعال کردن \lr{CoT}» کافی نیست. مدل‌های برتر (مثل \lr{Gemini}) یاد گرفته‌اند که «کوتاه و گزیده» استدلال کنند، در حالی که مدل‌های ضعیف‌تر در دام «نشخوار فکری» \lr{(Rumination)} و تولید متن‌های طولانی اما بیهوده می‌افتند. برای بهبود سیستم‌های حقوقی، باید مدل‌هایی را انتخاب کرد که تراکم اطلاعاتی \lr{(Information Density)} بالایی در استدلال‌های خود دارند و ستون‌های متوازنی در نمودار توکن دارند.

\end{document}